\phantomsection\section*{ПРИЛОЖЕНИЕ Б}\addcontentsline{toc}{section}{ПРИЛОЖЕНИЕ Б}

\begin{lstlisting}[language=C++, caption={Реализация алгоритма рендера части карта}, captionpos=b, label=list:render-task]
  for (int i = tile.i; i < tile.end_i; i++)
  {
    if (cancel_running)
      break;
    for (int j = tile.j; j < tile.end_j; j++)
    {
      if (cancel_running)
        break;
      Color pixel_color(0, 0, 0);
      for (int sample = 0; sample < _camera.samples_per_pixel; sample++)
      {
        if (cancel_running)
          break;
        Ray r = _camera.get_ray(j, i);
        pixel_color += _camera.ray_color(r, _camera.max_depth, world, cancel_running);
      }
      color_matrix.at(i, j) = _camera.pixel_samples_scale * pixel_color;
    }
\end{lstlisting}


\begin{lstlisting}[language=C++, caption={Реализация алгоритма получения луча из камеры}, captionpos=b, label=list:get-ray]
Ray get_ray(int i, int j) const
    {
        auto offset = sample_square();
        auto pixel_sample = pixel00_loc + ((i + offset.x()) * pixel_delta_u) + ((j + offset.y()) * pixel_delta_v);

        auto ray_origin = (_camera.defocus_angle <= 0) ? center : defocus_disk_sample();
        auto ray_direction = pixel_sample - ray_origin;

        return Ray(ray_origin, ray_direction);
    }
\end{lstlisting}

\begin{lstlisting}[language=C++, caption={Реализация алгоритма получения цвета  луча}, captionpos=b, label=list:ray-color]
Color ray_color(const Ray &r, int depth, const Hittable &world, volatile bool &cancel_running) const
    {
        if (cancel_running)
            return Color(0, 0, 0);

        if (depth <= 0)
            return Color(0, 0, 0);

        HitRecord rec;
        if (!world.hit(r, Interval(0.001, infinity), rec))
            return background;

        Ray scattered;
        Color attenuation;
        Color color_from_emission = rec.mat->emitted(rec.u, rec.v, rec.p);
        // Попытка рассеяния луча
        if (!rec.mat->scatter(r, rec, attenuation, scattered))
            return color_from_emission;

        Color color_from_scatter = attenuation * ray_color(scattered, depth - 1, world, cancel_running);

        return color_from_emission + color_from_scatter;
    }
\end{lstlisting}

\begin{lstlisting}[language=C++, caption={Реализация алгоритма нахождения пересечения с приземленным туманом}, captionpos=b, label=list:hit-fog]
Color ray_color(const Ray &r, int depth, const Hittable &world, volatile bool &cancel_running) const
  bool hit(const Ray &r, Interval ray_t, HitRecord &rec) const override
  {
    HitRecord rec1, rec2;
    if (!is_visible)
      return false;
    if (!boundary->hit(r, Interval::universe, rec1))
      return false;
    if (!boundary->hit(r, Interval(rec1.t + 0.0001, infinity), rec2))
      return false;

    auto t0 = std::fmax(rec1.t, ray_t.min);
    auto t1 = std::fmin(rec2.t, ray_t.max);
    if (t0 >= t1)
      return false;

    auto ray_len = r.direction().length();
    auto dist = (t1 - t0) * ray_len;

    auto hit_dist = -std::log(random_double()) / density;
    if (hit_dist > dist)
      return false;

    auto t = t0 + hit_dist / ray_len;
    Point3 p = r.at(t);

    double raw_noise = noise.noise(scale * p);
    double noise_val = 0.2 + 0.8 * std::fabs(raw_noise);
    double y = p.y();
    double height_factor = std::exp(-height_falloff * std::max(y, -10.0));

    double real_density = density * noise_val * height_factor;

    if (real_density < 1e-5 || random_double() > std::min(real_density, 1.0))
      return false;

    rec.t = t;
    rec.p = p;
    rec.normal = Vec3(0, 1, 0);
    rec.front_face = true;
    rec.mat = phase_function;
    return true;
}
\end{lstlisting}